\section{The heart commanded to love}

The command to love God with the whole heart stands inside a Mass that repeatedly asks God to heal the heart. The worshipper acknowledges divine justice, asks for mercy, hears the call to patience, confesses the Lord whom David himself worships, and approaches mysteries whose fruit is the cure of vice and an eternal remedy. Love is exacting: it reaches every part of the person and every neighbor. Yet the Mass does not leave the hearer alone with that demand. Its prayers ask for the freedom, purity, strength, and lasting healing that obedience requires.

The Gospel's closing question gives this movement its center. Jesus first answers what the law commands; then he asks who the Christ is. The answer ``David's son'' is true but incomplete when David, speaking in the Spirit, calls him Lord. The one who teaches love stands before his hearers as more than a teacher whose answer they can assess from a distance. His lordship claims their allegiance. The Epistle names one Lord within the common gifts that make the Church one; the Communion sets every earthly ruler beneath the sovereignty of God. The whole self and the whole community belong to the Lord they address.

Three readings unfold that belonging. The first follows the healing of divided desire: Augustine's account of grace, Chrysostom's exposition of the two love commands, and Thomas's Eucharistic teaching show how the commanded life is also a received gift. The second follows the gathering of a people: the divine Word and Spirit, the one body, and the common voice of supplication require humble mutual service. The third follows the worshipper's action: confession, petition, offering, reception, and continuing fidelity form one movement toward the eternal remedy. They answer different questions about the same Mass---how love becomes possible, whose common life it sustains, and how worship bears fruit after the celebration.

\begin{longtable}{>{\raggedright\arraybackslash}p{0.19\linewidth}>{\raggedright\arraybackslash}p{0.73\linewidth}}
\toprule
\textbf{Appointed element} & \textbf{Its movement}\\ \midrule
\endhead
Introit & Ps.~118:137,124; verse 1: justice confessed and mercy sought by a servant among those who walk God's law.\\
Collect & Freedom from diabolical contagion and a pure mind following God alone.\\
Epistle & Eph.~4:1--6: patient charity expresses the unity of body, Spirit, hope, Lord, faith, baptism, and Father.\\
Gradual & Ps.~32:12,6: an elected people and heavens established through the divine Word and Spirit.\\
Alleluia & Ps.~101:2: an afflicted voice asks that prayer reach God.\\
Gospel & Matt.~22:34--46: the twofold love command and the question of David's Son and Lord; Credo follows.\\
Offertory & Adapted Dan.~9:17--19: Daniel intercedes for the sanctuary and the people bearing God's name.\\
Secret & The holy action is asked to free the people from past and future offenses; Trinity Preface follows.\\
Communion & Ps.~75:12--13: vows paid and gifts brought to the Lord who humbles princes.\\
Postcommunion & Sacred mysteries heal vices and bring an eternal remedy.\\
\bottomrule
\end{longtable}

The Sunday is of the second class in the 1962 Missal. For the universal calendar on 20 September 2026 it is the Seventeenth Sunday after Pentecost; the ordinary commemoration of SS.~Eustace and Companions is omitted. The proper consists of the three prayers and seven scriptural elements above, whose texts follow with an explanation of each element and three sustained readings.\footnote{\textit{Missale Romanum} (Vatican, 1962), pp.~398--400, nos.~1602--1611; general rubrics, no.~111(b), p.~XVIII; September calendar, p.~LI. Dated corroboration: FSSP France and ICRSP France, Ordos for 20 September 2026, checked 17 September 2026.}

\tableofcontents
